% ============================================================================
% ARTIGO - SBSeg 2026 (Simpósio Brasileiro de Segurança da Informação)
% Formato: SBC (Sociedade Brasileira de Computação) - Coluna Única
% Limite: 14 páginas + 2 páginas para referências/apêndices = 16 total
% Revisão: Double-blind (anonimizar antes de submeter)
% Idioma: Português ou Inglês
% Site: https://www.sbseg2026.uff.br/chamadas/trilha-principal/
% ============================================================================

\documentclass[12pt]{article}

\usepackage{sbc-template}
\usepackage[utf8]{inputenc}
\usepackage[T1]{fontenc}
\usepackage[brazil]{babel}
\usepackage{graphicx}
\usepackage{url}
\usepackage{amsmath}
\usepackage{booktabs}
\usepackage{float}
\usepackage{enumitem}
\usepackage{placeins}
\usepackage{hyperref}
\usepackage{xcolor}

\sloppy

% ============================================================================
% TÍTULO
% ============================================================================
\title{Sistema Autônomo de Autenticação Contextual com Blockchain e\\Inteligência Artificial: Uma Abordagem Baseada em Microserviços\\para Segurança Adaptativa em Aplicações Web}

% ============================================================================
% AUTORES
% ATENÇÃO: Para submissão, REMOVER autores (double-blind).
% Descomente a linha abaixo para versão anônima:
% \author{}
% ============================================================================
\author{João Guedes de Brito\inst{1}, Cleber Jorge Lira de Santana\inst{1}}

\address{Programa de Pós-Graduação em Engenharia de Sistemas e Produtos\\
  Instituto Federal da Bahia (IFBA) -- Salvador, BA -- Brasil
  \email{\{joao.brito, cleber.santana\}@ifba.edu.br}
}

\begin{document}

\maketitle

% ============================================================================
% RESUMO (Português)
% ============================================================================
\begin{resumo}
Este artigo apresenta o desenvolvimento e a avaliação de um sistema autônomo de autenticação contextual que integra Inteligência Artificial (IA) e tecnologia Blockchain para reforçar a segurança de acesso em aplicações web. A arquitetura proposta baseia-se em microserviços e articula cinco camadas funcionais: aplicação modular, integração contínua, segurança automatizada, inteligência artificial e registro imutável. O módulo de IA emprega um modelo \textit{ensemble} (Random Forest, Decision Tree e Logistic Regression) para classificação de risco em tempo real, analisando dez atributos contextuais --- incluindo geolocalização, dispositivo, horário e padrões comportamentais. A camada Blockchain garante imutabilidade e rastreabilidade forense dos logs de autenticação, atendendo a requisitos normativos como a LGPD. Os resultados experimentais demonstram acurácia de 91\% na detecção de acessos anômalos, latência média de 320ms no fluxo completo de autenticação e integridade de 100\% nas transações Blockchain. A solução demonstra viabilidade técnica como alternativa adaptativa e auditável aos métodos tradicionais de autenticação.
\end{resumo}

% ============================================================================
% ABSTRACT (Inglês)
% ============================================================================
\begin{abstract}
This paper presents the development and evaluation of an autonomous contextual authentication system integrating Artificial Intelligence (AI) and Blockchain technology to enhance access security in web applications. The proposed architecture is based on microservices and comprises five functional layers: modular application, continuous integration, automated security, artificial intelligence, and immutable logging. The AI module employs an ensemble model (Random Forest, Decision Tree, and Logistic Regression) for real-time risk classification, analyzing ten contextual attributes --- including geolocation, device, time, and user behavioral patterns. The Blockchain layer ensures immutability and forensic traceability of authentication logs, meeting regulatory requirements such as Brazil's LGPD. Experimental results demonstrate 91\% accuracy in anomalous access detection, with an average latency of 320ms and 100\% integrity in Blockchain-recorded transactions.
\end{abstract}

% ============================================================================
% 1. INTRODUÇÃO
% ============================================================================
\section{Introdução}

Em um cenário digital marcado pela crescente sofisticação dos ataques cibernéticos e pela expansão dos serviços baseados em web, a autenticação de usuários permanece como um dos pilares mais críticos da segurança da informação. Métodos tradicionais, como senhas fixas e autenticação em dois fatores (2FA), embora amplamente adotados, apresentam limitações significativas diante de ataques como \textit{credential stuffing}, \textit{session hijacking} e engenharia social \cite{owasp2023}. Essas abordagens, predominantemente estáticas, não consideram o contexto comportamental e ambiental do usuário no momento do acesso, criando lacunas exploráveis por agentes maliciosos.

A crescente demanda por conformidade regulatória, como a Lei Geral de Proteção de Dados (LGPD) no Brasil \cite{brasil2018}, e as diretrizes internacionais do National Institute of Standards and Technology (NIST) \cite{nist2020}, reforçam a necessidade de sistemas de autenticação que combinem adaptabilidade, auditabilidade e transparência. Nesse contexto, a convergência entre Inteligência Artificial (IA), tecnologia Blockchain e arquiteturas distribuídas emerge como abordagem promissora para superar as limitações dos mecanismos convencionais.

Este trabalho propõe e implementa um sistema autônomo de autenticação contextual que integra: (i) análise de risco em tempo real baseada em modelos supervisionados de Machine Learning; (ii) registro imutável de eventos de autenticação via Blockchain; e (iii) uma arquitetura de microserviços orientada por práticas DevSecOps. O sistema analisa dez atributos contextuais por acesso --- incluindo localização geográfica, dispositivo, horário, padrão comportamental e indicadores de rede --- para tomar decisões autônomas sobre a concessão, monitoramento, exigência de autenticação multifator ou bloqueio de acesso.

A principal contribuição deste trabalho reside na integração sistêmica de três camadas complementares de segurança --- autenticação baseada em identidade (Keycloak/OAuth2), análise preditiva de risco (IA) e auditoria imutável (Blockchain) --- em uma arquitetura modular, escalável e reprodutível. Os resultados experimentais demonstram a viabilidade técnica da abordagem, com acurácia de 91\% na classificação de acessos anômalos e latência compatível com ambientes de produção.

O restante deste artigo está organizado da seguinte forma: a Seção 2 apresenta trabalhos relacionados; a Seção 3 descreve a fundamentação teórica; a Seção 4 detalha a arquitetura e implementação; a Seção 5 apresenta a metodologia e os resultados experimentais; e a Seção 6 conclui o trabalho.

% ============================================================================
% 2. TRABALHOS RELACIONADOS
% ============================================================================
\section{Trabalhos Relacionados}

A autenticação contextual e adaptativa tem sido objeto de investigação crescente na literatura de segurança da informação. Abordagens baseadas em análise comportamental para detecção de anomalias em processos de autenticação são discutidas por \cite{goodfellow2016}, que demonstram a eficácia de modelos supervisionados na captura de padrões não lineares em dados de alta dimensionalidade. Entretanto, a maioria dessas abordagens trata a análise de risco de forma isolada, sem integração com mecanismos de auditoria imutável.

No domínio da Blockchain aplicada à segurança, o relatório técnico NISTIR 8202 \cite{nist2018blockchain} descreve a aplicabilidade de blockchains permissionadas em ambientes corporativos para registro de evidências e trilhas de auditoria. Trabalhos recentes exploram o uso de Blockchain para registro de logs de autenticação, porém tipicamente não incorporam análise preditiva de risco baseada em IA como camada complementar de decisão.

Soluções comerciais como o Azure AD Identity Protection (Microsoft) e o Google BeyondCorp implementam análise contextual parcial, porém operam como sistemas proprietários e centralizados, sem oferecer transparência nos critérios decisórios nem auditabilidade descentralizada. O \textit{AI Risk Management Framework} do NIST \cite{nist2023ai} destaca a importância da rastreabilidade e governança em sistemas baseados em IA, requisitos que sistemas proprietários frequentemente não atendem.

O presente trabalho diferencia-se ao propor a integração sistêmica de três camadas de segurança --- identidade (Keycloak/OAuth2), inteligência preditiva (\textit{ensemble} de ML) e auditoria imutável (Blockchain) --- em uma arquitetura aberta, modular e reprodutível, com decisões automatizadas baseadas em critérios normativos, estatísticos e criptográficos.

% ============================================================================
% 3. FUNDAMENTAÇÃO TEÓRICA
% ============================================================================
\section{Fundamentação Teórica}

\subsection{Autenticação Contextual e Segurança Adaptativa}

A autenticação contextual representa uma evolução dos modelos tradicionais ao incorporar variáveis ambientais e comportamentais no processo decisório de concessão de acesso. O \textit{Secure Software Development Framework} (SSDF) do NIST \cite{nist2022} estabelece diretrizes para desenvolvimento seguro que incluem automação de controles e integração de mecanismos de verificação contínua. A abordagem \textit{Secure by Design}, incentivada pela CISA \cite{cisa2023}, propõe que a segurança seja considerada desde a concepção do sistema.

\subsection{Machine Learning para Detecção de Anomalias}

Técnicas de Machine Learning permitem identificar padrões anômalos em comportamento de usuários, aumentando a capacidade preditiva dos sistemas de defesa \cite{nist2023ai}. Modelos supervisionados classificam vulnerabilidades conhecidas, enquanto modelos não supervisionados possibilitam detecção de comportamentos anômalos não catalogados. Modelos \textit{ensemble}, que combinam múltiplos classificadores por mecanismos de votação, mitigam vieses individuais e aumentam a robustez preditiva \cite{scikit2024}.

\subsection{Blockchain para Integridade de Logs}

A tecnologia Blockchain, fundamentada no encadeamento criptográfico descrito por Nakamoto \cite{nakamoto2008}, oferece garantias de imutabilidade e não repúdio para registros digitais. Blockchains permissionadas são adequadas para ambientes corporativos, combinando controle de acesso institucional com propriedades de integridade distribuída \cite{nist2018blockchain}. A utilização de Blockchain para registrar \textit{hashes} de eventos de autenticação garante que nenhum registro seja alterado retroativamente.

\subsection{DevSecOps e Microserviços}

O DevSecOps integra controles de segurança desde as fases iniciais do desenvolvimento, promovendo \textit{Shift Left Security} \cite{nist2020}. A incorporação de análise estática (SAST), dinâmica (DAST) e escaneamento de dependências reduz a exposição a vulnerabilidades do OWASP Top 10 \cite{owasp2023}. A arquitetura de microserviços \cite{newman2021} favorece decomposição funcional, isolamento de domínios e aplicação de políticas granulares de segurança.

% ============================================================================
% 4. ARQUITETURA E IMPLEMENTAÇÃO
% ============================================================================
\section{Arquitetura e Implementação}

\subsection{Visão Geral da Arquitetura em Cinco Camadas}

A arquitetura proposta organiza-se em cinco camadas integradas:

\begin{enumerate}[leftmargin=*]
    \item \textbf{Camada de Aplicação Modular} --- serviços independentes (backend Spring Boot 3.2.3, frontend Angular 17) e APIs autenticadas via JWT/OAuth2, com Keycloak 24 como provedor de identidade.
    \item \textbf{Camada de Integração Contínua} --- orquestração automatizada do pipeline DevSecOps com testes unitários, de integração e de segurança.
    \item \textbf{Camada de Segurança Automatizada} --- verificações SAST, DAST e análise de dependências contra bases de vulnerabilidades conhecidas.
    \item \textbf{Camada de Inteligência Artificial} --- classificação de riscos e detecção de anomalias com modelos supervisionados de ML.
    \item \textbf{Camada de Registro Imutável} --- rastreabilidade e integridade forense via Blockchain permissionada (Ethereum/Hyperledger Fabric).
\end{enumerate}

% Descomente para incluir a figura da arquitetura:
% \begin{figure}[H]
%   \centering
%   \includegraphics[width=0.95\textwidth]{figuras/arquitetura-cinco-camadas.png}
%   \caption{Arquitetura em Cinco Camadas do Sistema Proposto}
%   \label{fig:arquitetura}
% \end{figure}

O fluxo operacional segue uma cadeia decisória automatizada: a tentativa de autenticação ativa a captura de contexto; os atributos são processados pelas estratégias de análise; o modelo de IA calcula a probabilidade de risco; critérios objetivos de conformidade são aplicados; o resultado é registrado na Blockchain; e o sistema decide pela liberação, monitoramento, MFA ou bloqueio do acesso.

\subsection{Módulo de Autenticação Contextual}

O módulo de autenticação opera como intermediário inteligente no fluxo de login. O sistema captura automaticamente dez atributos contextuais por tentativa de acesso, conforme a Tabela~\ref{tab:features}.

\begin{table}[H]
\centering
\caption{Atributos contextuais extraídos por tentativa de autenticação}
\label{tab:features}
\small
\begin{tabular}{@{}clp{7cm}@{}}
\toprule
\textbf{\#} & \textbf{Atributo} & \textbf{Descrição} \\
\midrule
1  & \texttt{hour\_of\_day}              & Hora do acesso (0--23) \\
2  & \texttt{day\_of\_week}              & Dia da semana (0--6) \\
3  & \texttt{login\_attempts\_last\_hour} & Tentativas de login na última hora \\
4  & \texttt{is\_new\_device}            & Dispositivo não reconhecido (0/1) \\
5  & \texttt{is\_new\_location}          & Localização inédita para o usuário (0/1) \\
6  & \texttt{is\_vpn}                    & Uso de VPN detectado (0/1) \\
7  & \texttt{is\_tor}                    & Uso de rede Tor detectado (0/1) \\
8  & \texttt{session\_duration\_avg}     & Duração média das sessões anteriores (s) \\
9  & \texttt{pages\_per\_session\_avg}   & Páginas por sessão (média histórica) \\
10 & \texttt{time\_since\_last\_login}   & Tempo desde o último login (horas) \\
\bottomrule
\end{tabular}
\end{table}

O backend implementa cinco estratégias de análise de risco seguindo o padrão \textit{Strategy}: análise comportamental, de dispositivo, de localização, de rede e temporal. Cada estratégia contribui com uma pontuação parcial que alimenta o modelo de IA para classificação final.

\subsection{Modelo de Machine Learning}

O serviço de IA foi implementado como microsserviço independente em Python com FastAPI 0.109.0, expondo endpoints RESTful para predição (\texttt{/predict}), treinamento (\texttt{/train}) e verificação de saúde (\texttt{/health}).

O modelo utiliza \textit{ensemble} baseado em \textit{Soft Voting}, combinando três algoritmos:

\begin{itemize}[leftmargin=*]
    \item \textbf{Random Forest} --- 100 estimadores, profundidade máxima 10;
    \item \textbf{Decision Tree} --- profundidade máxima 8;
    \item \textbf{Logistic Regression} --- 500 iterações, otimizador \textit{lbfgs}.
\end{itemize}

O mecanismo de \textit{Soft Voting} calcula a média das probabilidades preditas por cada classificador, resultando em um \textit{score} de risco contínuo $s \in [0, 1]$. Os dados são normalizados via \texttt{StandardScaler} antes da inferência. O modelo foi treinado com 5.000 amostras sintéticas geradas a partir de regras heurísticas que simulam cenários realistas.

A classificação segue quatro níveis de risco (Tabela~\ref{tab:risk_levels}), cada um associado a uma decisão automatizada:

\begin{table}[H]
\centering
\caption{Níveis de risco e decisões automatizadas do sistema}
\label{tab:risk_levels}
\small
\begin{tabular}{@{}llll@{}}
\toprule
\textbf{Nível} & \textbf{Score ($s$)} & \textbf{Decisão} & \textbf{Ação} \\
\midrule
Baixo    & $s < 0{,}3$          & PERMITIR                   & Acesso liberado \\
Médio    & $0{,}3 \leq s < 0{,}6$ & MONITORAR                  & Acesso com alerta \\
Alto     & $0{,}6 \leq s < 0{,}8$ & EXIGIR\_MFA                & Autenticação multifator \\
Crítico  & $s \geq 0{,}8$       & BLOQUEAR                   & Acesso negado \\
\bottomrule
\end{tabular}
\end{table}

Fatores de risco identificados automaticamente incluem: uso de Tor/VPN, dispositivo desconhecido, localização inédita, múltiplas tentativas de login, horário atípico (2h--6h) e inatividade prolongada ($>$30 dias).

\subsection{Camada Blockchain}

A camada de registro imutável suporta duas implementações:

\begin{itemize}[leftmargin=*]
    \item \textbf{Ethereum} (via Web3j 4.10.3) --- registro de \textit{hashes} criptográficos em rede pública ou privada;
    \item \textbf{Hyperledger Fabric 2.5} --- governança permissionada para ambientes corporativos.
\end{itemize}

Cada evento de autenticação gera uma transação contendo: tipo do evento, identificador do usuário (anonimizado via \textit{hash} para conformidade com LGPD), endereço IP, \textit{fingerprint} do dispositivo, \textit{score} de risco, decisão tomada e \textit{timestamp}. A integridade é verificável por endpoint que realiza validação criptográfica da cadeia, com confirmação mínima de 12 blocos.

\subsection{Gerenciamento de Identidade e Observabilidade}

O Keycloak 24 foi configurado como provedor central de identidade, suportando OAuth 2.0 e OpenID Connect, com proteção contra força bruta (5 tentativas, bloqueio de 900s), PKCE para o frontend e tokens de curta duração (5 minutos).

O monitoramento integra ELK Stack 8.12 (Elasticsearch, Logstash, Kibana) e Grafana 10.3 para observabilidade em tempo real. O pipeline Logstash classifica eventos por tipo (autenticação, risco, blockchain, alertas) e indexa no Elasticsearch com padrão temporal.

\subsection{Stack Tecnológica}

A Tabela~\ref{tab:tech_stack} resume as tecnologias empregadas.

\begin{table}[H]
\centering
\caption{Stack tecnológica do sistema proposto}
\label{tab:tech_stack}
\small
\begin{tabular}{@{}lll@{}}
\toprule
\textbf{Camada} & \textbf{Tecnologia} & \textbf{Versão} \\
\midrule
Backend           & Spring Boot        & 3.2.3 \\
Frontend          & Angular            & 17.2.0 \\
Identidade        & Keycloak           & 24.0 \\
ML/IA             & Scikit-learn       & 1.4.0 \\
API de IA         & FastAPI            & 0.109.0 \\
Blockchain        & Web3j (Ethereum)   & 4.10.3 \\
Blockchain        & Hyperledger Fabric & 2.5 \\
Banco de Dados    & PostgreSQL         & 15 \\
Cache/Sessões     & Redis              & 7 \\
Containerização   & Docker Compose     & Latest \\
Logs              & ELK Stack          & 8.12.0 \\
Monitoramento     & Grafana            & 10.3.1 \\
\bottomrule
\end{tabular}
\end{table}

% ============================================================================
% 5. AVALIAÇÃO EXPERIMENTAL
% ============================================================================
\section{Avaliação Experimental}

\subsection{Configuração do Ambiente}

O ambiente experimental foi construído com infraestrutura containerizada via Docker, com dez microsserviços orquestrados em rede bridge compartilhada: PostgreSQL, Redis, Keycloak, backend Spring Boot, frontend Angular, serviço de IA (FastAPI), Elasticsearch, Logstash, Kibana e Grafana.

\subsection{Metodologia}

Os experimentos foram estruturados em três eixos de validação:

\begin{enumerate}[leftmargin=*]
    \item \textbf{Avaliação do modelo de ML}: classificação de acessos legítimos e anômalos com validação cruzada ($k$-\textit{fold}, $k=10$);
    \item \textbf{Cadeia de autenticação contextual}: integração completa Keycloak $\rightarrow$ captura de contexto $\rightarrow$ classificação por IA $\rightarrow$ decisão $\rightarrow$ registro Blockchain;
    \item \textbf{Integridade Blockchain}: verificação de imutabilidade e consistência criptográfica das transações.
\end{enumerate}

O modelo foi treinado com 5.000 amostras sintéticas contendo os dez atributos contextuais, rotuladas em quatro categorias de risco com base em regras derivadas de cenários reais de ataque. Foram simulados cenários controlados: acesso legítimo habitual; acesso de novo dispositivo em localização conhecida; acesso via VPN/Tor em horário atípico; múltiplas tentativas falhas (\textit{credential stuffing}); e acesso após inatividade prolongada.

\subsection{Resultados}

\subsubsection{Desempenho do Modelo de ML}

Os resultados do modelo \textit{ensemble} são apresentados na Tabela~\ref{tab:ml_results}.

\begin{table}[H]
\centering
\caption{Métricas de desempenho do modelo de classificação}
\label{tab:ml_results}
\small
\begin{tabular}{@{}ll@{}}
\toprule
\textbf{Métrica} & \textbf{Valor} \\
\midrule
Acurácia            & 91,0\% \\
Precisão            & 89,4\% \\
Recall              & 93,1\% \\
F1-Score            & 91,2\% \\
Taxa de Falsos Positivos & 8,6\% \\
\bottomrule
\end{tabular}
\end{table}

O F1-\textit{score} de 91,2\% indica equilíbrio entre precisão e sensibilidade. O \textit{recall} de 93,1\% é particularmente relevante em segurança, pois prioriza a detecção de acessos genuinamente anômalos, minimizando falsos negativos. A taxa de falsos positivos de 8,6\% é aceitável para validação com dados sintéticos e tende a reduzir com dados reais.

\subsubsection{Latência da Autenticação Contextual}

A latência média do fluxo completo foi de 320ms, conforme decomposição na Tabela~\ref{tab:latency}.

\begin{table}[H]
\centering
\caption{Decomposição da latência do fluxo de autenticação}
\label{tab:latency}
\small
\begin{tabular}{@{}ll@{}}
\toprule
\textbf{Etapa} & \textbf{Latência Média} \\
\midrule
Validação de credenciais (Keycloak)      & 80--100ms \\
Captura de contexto                       & 20--30ms \\
Classificação de risco (IA)               & 120--180ms \\
Registro em Blockchain                    & 40--60ms \\
\midrule
\textbf{Total}                            & \textbf{$\approx$320ms} \\
\bottomrule
\end{tabular}
\end{table}

O módulo de IA (120--180ms) é o principal gargalo. Estratégias de mitigação incluem cache de decisões para padrões reconhecidos e quantização do modelo. O tempo total de 320ms é compatível com padrões de UX em autenticação web (tolerância de até 500ms).

\subsubsection{Integridade Blockchain}

Os testes demonstraram consistência de 100\% nas transações registradas, com validação criptográfica bem-sucedida em todos os registros. A anonimização por \textit{hash} dos dados pessoais antes do registro atende ao princípio de minimização da LGPD.

\subsection{Análise Comparativa}

A Tabela~\ref{tab:comparison} compara o sistema proposto com abordagens tradicionais.

\begin{table}[H]
\centering
\caption{Comparação entre abordagens de autenticação}
\label{tab:comparison}
\small
\begin{tabular}{@{}lccc@{}}
\toprule
\textbf{Característica} & \textbf{Senha} & \textbf{2FA} & \textbf{Proposta} \\
\midrule
Análise contextual       & Não & Não & Sim \\
Adaptabilidade           & Não & Não & Sim \\
Auditoria imutável       & Não & Não & Sim \\
Detecção de anomalias    & Não & Parcial & Sim \\
Conformidade LGPD        & Parcial & Parcial & Sim \\
Resistência a \textit{credential stuffing} & Não & Parcial & Sim \\
Decisão autônoma         & Não & Não & Sim \\
\bottomrule
\end{tabular}
\end{table}

\subsection{Discussão}

Os resultados confirmam que a integração de IA contextual e Blockchain ao processo de autenticação eleva significativamente o nível de segurança em comparação com métodos estáticos. A abordagem adaptativa identificou padrões suspeitos em cenários de \textit{credential stuffing} e \textit{session hijacking} que seriam invisíveis a sistemas convencionais.

A arquitetura de microserviços demonstrou adequação para escalabilidade independente dos componentes. O Keycloak ofereceu suporte nativo a OAuth2 e OpenID Connect sem desenvolvimento adicional. A camada Blockchain demonstrou robustez para rastreabilidade forense, garantindo que nenhum registro possa ser adulterado retroativamente.

As limitações identificadas incluem: (i) utilização de dados sintéticos para treinamento, impactando a generalização; (ii) complexidade operacional do Hyperledger Fabric; e (iii) necessidade de políticas refinadas de anonimização para plena conformidade com a LGPD.

% ============================================================================
% 6. CONCLUSÃO E TRABALHOS FUTUROS
% ============================================================================
\section{Conclusão e Trabalhos Futuros}

Este trabalho apresentou o desenvolvimento e a avaliação de um sistema autônomo de autenticação contextual integrando IA e Blockchain em arquitetura de microserviços com práticas DevSecOps. Os resultados demonstram viabilidade técnica com acurácia de 91\%, latência de 320ms compatível com produção e integridade total dos registros Blockchain.

A contribuição principal reside na integração sistêmica de três camadas complementares --- identidade, inteligência preditiva e auditoria imutável --- em sistema modular e reprodutível, oferecendo um paradigma proativo e orientado por evidências para autenticação em ambientes críticos.

Como trabalhos futuros: (i) coleta de dados reais para retreinamento do modelo; (ii) incorporação de aprendizado contínuo (\textit{online learning}); (iii) testes de carga em alta concorrência; (iv) extensão para IoT e autenticação federada; e (v) implementação de explicabilidade (XAI) para as decisões do modelo.

% ============================================================================
% REFERÊNCIAS
% ============================================================================
\bibliographystyle{sbc}

\begin{thebibliography}{99}

\bibitem{brasil2018}
BRASIL. Lei nº 13.709, de 14 de agosto de 2018. Lei Geral de Proteção de Dados Pessoais (LGPD). \textit{Diário Oficial da União}, Brasília, DF, 2018.

\bibitem{cisa2023}
CISA. Secure by Design Initiative. Washington, DC: CISA, 2023.

\bibitem{eua2025}
ESTADOS UNIDOS. Executive Order 14144: Strengthening and Promoting Innovation in the Nation's Cybersecurity. Washington, DC: The White House, 2025.

\bibitem{goodfellow2016}
Goodfellow, I., Bengio, Y., and Courville, A. (2016). \textit{Deep Learning}. MIT Press, Cambridge.

\bibitem{hyperledger2024}
Hyperledger Foundation. (2024). Hyperledger Fabric Documentation, v2.5. Linux Foundation.

\bibitem{keycloak2024}
Keycloak. (2024). Server Administration Guide, v24. Red Hat.

\bibitem{nakamoto2008}
Nakamoto, S. (2008). Bitcoin: A Peer-to-Peer Electronic Cash System.

\bibitem{nist2023ai}
NIST. (2023). AI Risk Management Framework (AI RMF 1.0). NIST AI 100-1.

\bibitem{nist2018blockchain}
NIST. (2018). Blockchain Technology Overview. NISTIR 8202.

\bibitem{nist2022}
NIST. (2022). Secure Software Development Framework (SSDF) v1.1. SP 800-218.

\bibitem{nist2020}
NIST. (2020). Security and Privacy Controls for Information Systems and Organizations. SP 800-53 Rev. 5.

\bibitem{newman2021}
Newman, S. (2021). \textit{Building Microservices: Designing Fine-Grained Systems}. 2nd ed. O'Reilly Media.

\bibitem{owasp2023}
OWASP Foundation. (2023). OWASP Top 10 -- 2023 Edition.

\bibitem{scikit2024}
Scikit-learn Developers. (2024). Scikit-learn: Machine Learning in Python, v1.4.

\end{thebibliography}

\end{document}
